
\documentclass{article}

\usepackage{polski}
\usepackage[utf8]{inputenc}
\usepackage{geometry}
\usepackage{graphicx}
\usepackage{amssymb}
\usepackage{amsmath}
\usepackage{amsthm}
\usepackage{empheq}
\usepackage{mdframed}
\usepackage{booktabs}
\usepackage{lipsum}
\usepackage{graphicx}
\usepackage{color}
\usepackage{psfrag}
\usepackage{pgfplots}
\usepackage{bm}
\usepackage{array, float}
\usepackage[style=authoryear, backend=biber]{biblatex}
%%%%%%%%%%%%%%%%%%%%%%%%%%%%%%%%%%%%%%%%%%%%%%%%%%%%%%%%%%%%%%%%%%%%%%%%%%%%%%%

% Other Settings

%%%%%%%%%%%%%%%%%%%%%%%%%% Page Setting %%%%%%%%%%%%%%%%%%%%%%%%%%%%%%%%%%%%%%%
\geometry{a4paper}
\addbibresource{bibliografia.bib}
%%%%%%%%%%%%%%%%%%%%%%%%%% Define an orangebox command %%%%%%%%%%%%%%%%%%%%%%%%
\newcommand\orangebox[1]{\fcolorbox{ocre}{mygray}{\hspace{1em}#1\hspace{1em}}}
%%%%%%%%%%%%%%%%%%%%%%%%%%%%%%%%%%%%%%%%%%%%%%%%%%%%%%%%%%%%%%%%%%%%%%%%%%%%%%%


\title{Porównanie efektywności algorytmów sortowania}
\author{Krystian Mikołajczak 284452}
%%%%%%%%%%%%%%%%%%%%%%%%%%%%%%%%%%%%%%%%%%%%%%%%%%%%%%%%%%%%%%%%%%%%%%%%%%%%%%%

\begin{document}
    \maketitle
  \section{Wprowadzenie}
      \subsection{Opis zadania}
        Celem jest zaimplementowanie wybranych algorytmów sortowania i przeprowadzenie analizy ich efektywności.
        Uzyskane w wyniku badań realne czasy sortowania dla wybranych algorytmów zostaną poddane analizie oraz porównane wyników innych algorymów.
        Zbadany zostanie także wpływ typu danych na efektywność jednego wybranego algorytmu.
    \subsection{Algorytmy poddane analizie}
      Otrzymane wyniki zostały porównane z teoretyczną złożonością obliczeniową analizowanych algorytmów. Analizie został także podany wpływ danych poddanych sortowaniu na efektywność algorytmów.
      \subsubsection{Sortowanie przez kopcowanie} 
        Algorytm sortowanie przez kopcowanie możemy podzielić na dwie części, w pierwszej kolejności musimy sortowaną tablicę doprowadzić do poprawnej struktury kopca maksymalnego,
        czyli takiej w której każdy węzeł jest mniejszy lub równy swojemu rodzicowi.
        Następnie z utworzonej struktury kopca usuwamy korzeń, będący największym elementem w kopcu po czym zamieniamy go z ostatnim elementem w kopcu i zmniejszamy rozmiar kopca. Po każdym wykonaniu takiej operacji musimy przywrócić własność kopca zaczynając od nowego korzenia. 
        \\\\Algorytm ten, składając w całość tworzenie kopca o złożoności $O(n)$ 
        wywoływane raz oraz przywracanie własności kopca o złożności
        $O(\log n)$ wywoływane $n$ razy, daje nam stałą teoretyczną
        wydajność dla średniego i najgorszego przypadku w których złożoność
        obliczeniowa jest równa $O(n \log n)$ [\cite{AverageHeapSort}] [\cite{WorstHeapSort}] 
        \\\\Zakłada się więc, że czas sortowania dla tego algorytmu będzie stabilnie wzrastał wraz ze wzrostem wielkości badanej instancji, wykazując niewrażliwość na stopień wstępnego posortowania tablicy.
      \subsubsection{Sortowanie przez wstawianie}
        Algorytm sortowania przez wstawianie buduje posortowaną tablicę element po elemencie. Algorytm iteruje przez całą sortowaną tablicę, każdy element jest usuwany z części tablicy do posortowania i wstawiany we właściwe miejsce w posortowanej częsci tablicy.
        Efekt ten jest osiągany poprzez porównywanie aktualnie przenoszonego elemetnu z każdym elementem znajdującym się w posortowanej już części tablicy aż do momentu gdy nie zostanie napotkany element mniejszy niż aktualnie przenoszzony, który zostaje wstawiony w puste miejsce.
        \\\\Rezultatem tego działania jest teoretyczna złożoność obliczeniowa w średnim i najgorszym przypadku na poziomie $O(n^2)$,
        zaletą tego algorytmu jest jednak wydajność teoretyczna w przypadku prawie posortowanej tablicy na poziomie $O(n)$. [\cite{wikiInsert}]
        \\\\Zakłada się więc, że wystąpią drastyczne różnice w wydajności tego algorytmu w zależności od wstępnego stopnia posortowania tablicy, dla posortowanej tablicy oczekiwana jest złożoność na poziomie $O(n)$, a dla tablicy odwrotnie posortowanej oczekiwana jest złożoność rzędu $O(n^2)$
      \subsubsection{Sortowanie szybkie}
        Sortowanie szybkie jest algorytmem działającym na zasadzie "dziel i zwyciężaj",
        polegającej na rekurencyjnym dzieleniu problemu na mniejsze podproblemy tego samego typu tak długo,
        aż podproblemy te staną się wystarczająco proste do rozwiązania.
        W przeciwieństwie do innych algorytmów tego typu, w sortowaniu szybkim głowne porządkowanie odbywa się w trakcie tego podziału, a nie w trakcie późniejszego scalania rozwiązań podproblemów. [\cite{wikiDziel}]
        \\\\Z sortowanej tablicy wybierany jest element który dzieli ją na dwa fragmenty nazywany pivotem lub elementem osiowym. Do jednego przenoszone są wszystkie elementy mniejsze od wybranego elementu a do drugiego te większe.Po takim podziale element osiowy znajduje się na właściwym miejscu. 
        Powoduje to, że wydajność tego algorytmu jest w dużej mierze zależna od wyboru tego właśnie elementu.
        Następnie schemat ten powtarzany jest dla każdego z dwóch otrzymanych fragmentów. Powtarzanie dzielenia kończy się w momencie gdy otrzymana z podziału część posiada tylko jeden element.
        \\\\Algorytm ten, biorąc pod uwagę jego rekursywną naturę w teoretycznym najgorszym przypadku przy złym wyborze elementu osiowego daje nam wydajność na poziomie $O(n^2)$, za to teoretyczna średnia złożoność jest na poziomie $O(n\log{n})$ 
        \\\\Zakłada się więc, że sortowanie szybkie osiągnie najlepsze czasy dla danych losowych osiągając złożoność na poziomie $n\log{n}$, jednak przy złym doborze elementu osiowego spodziewany jest spadek nawet do okolic $O(n^2)$ 
      \subsubsection{Sortowanie Shella - wyrażenie $2^k - 1$, ciąg opracowany przez Marcina Ciurę}
        Algorytm sortowania Shella korzysta z porównań, działa podobnie do algorytmu sortowania przez wstawianie, jednak kluczową różnicą jest to że w algorytmie Shella porównujemy i zamieniamy najpierw elementy położone daleko od siebie, stopniowo zmniejszając tę odległość po każdym przejściu przez sortowaną tablicę.
        \\\\Wydajność tego algorytmu zależy w głównej mierze od doboru sekwencji odstępów 
        między elementami na których będziemy operować. W tym przypadku analizie zostaną
        poddane dwa ciągi, pierwszy o wyrażeniu ogólnym
        $2^k - 1$ opracowany przez Thomasa Hibbarda w 1963 r. [\cite{10.1145/366552.366557}],
        oraz drugi wyznaczony przez Marcina Ciurę w 2001 r.
        w postaci ${1,4,10,23,57,132,301,701}$ [\cite{CiuraShell}]
        \\\\
        W przypadku wyrażenia ogólnego  $2^k - 1$ teoretyczna pesymistyczna złożoność plasuje się na poziomie $O(n^{3/2})$ a złożoność średnia, zbadana doświadczalnie przez Thomasa Hibbarda, na poziomie $O(n^{1.25})$.
        Dla ciągu opracowanego przez Marcina Ciurę nie jesteśmy jednak w stanie podać tak dokładnych oczekiwań, dla tablic o rozmiarze do 128 tys. elementów średnia wydajność osiąga wartości w okolicach od $O(n^{1.1})$ do $O(n\log^2{n})$ a szacowana pesymistyczna złożoność wynosi około $O(n^{1.5})$
        Ciąg opracowany przez Marcina Ciurę jest o tyle wyjątkowy że został opracowany empirycznie, Autor porównywał wyniki złożoności sortowania dla małych tablic w poszukiwaniu najlepszego wyniku dla przypadku średniego.
        \\\\Zakłada się więc, że wydajność sortowania tym Algorytmem będzie znacznie wyższa niż przy sortowaniu przez wstawianie, oraz że zastosowanie ciągu Marcina Ciury wykaże lepsze wyniki niż sekwencja Hibbarda.
    \subsection{Charakterystyka danych badawczych}
      Dane wykorzystane w badaniu zostały wygenerowane przy użyciu generatora liczb pseudolosowych. Podstawowym badanym typem danych jest format liczb całkowitych (int) zajmujący 4 bajty pamięci,
      dla algorytmu sortowania szybkiego przeprowadzono także badanie stosując liczby zmiennoprzecinkowe (float, 4 bajty), liczby zmiennoprzecinkowe podwójnej precyzji (double, 8 bajtów), oraz znaki ASCII (char, 1 bajt).
      Badanie każdego algorytmu zostało przeprowadzone na tych samych instancjach tablic aby mieć możliwość porównania wyników między badanymi algorytmami.
      \\\\W badaniu uwzględniono pięć rodzajów rozkładu danych w tablicach:
      \begin{itemize}
        \item Dane losowe - zakres od 0 do 2,147,483,647
        \item Dane posortowane w 33\% - posrotowana część wypełniona zerami, reszta tablicy losowa w zakresie od 1 do 2,147,483,647
        \item Dane posortowane w 66\% - posrotowana część wypełniona zerami, reszta tablicy losowa w zakresie od 1 do 2,147,483,647
        \item Dane posortowane rosnąco - tablica wypełniona liczbami od 1 do rozmiaru tablicy
        \item Dane posortowane malejąco - tablica wypełniona liczbami od rozmiaru tablicy do 1
      \end{itemize}  
      Do badania wybrano siedem instancji o rozmiarach od $5 \cdot 10^3$ do $5 \cdot 10^4$ elementów. Wybór ten został podyktowany faktem że dolny próg jest wystarczająco duży aby można zmierzyć dokładnie czas sortowania, a górny próg pozwala na wyraźne zaobserwowanie różnicy między algorytmami, co pozwala na weryfikację postawionych hipotez.
      Taki dobór wielkości instancji pozwala jednocześnie uniknąć problemów związanych z przepełnieniem stosu na wykorzystanym do przeprowadzenia badań sprzęcie.
  \section{Warunki badawcze}
    \subsection{Środowisko badawcze}
      Badanie zostało przeprowadzone za pomocą autorskiego oprogramowania napisanego w języku C++, skompilowanego za pomocą kompilatora Apple Clang w wersji 17.0.0 ,zgodnie ze standardem C++20, stosując standardową konfigurację.
      Specyfikacja platformy badawczej:
      \begin{itemize}
        \item \textbf{Procesor:} Apple M4, 10 rdzeni (4 rdzenie wydajnościowe, 6 rdzeni energooszczędnych)
      \item \textbf{Pamięć operacyjna:} 16 GB zunifikowanej pamięci RAM
      \item \textbf{System operacyjny:} macOS Tahoe 26.3.1
        \item \textbf{Chłodzenie:} brak aktywnego chłodzenia - w czasie testów monitorowano temperaturę za pomocą zewnętrznego oprogramowania aby uniknąć ograniczenia wydajności z powodu zbyt wysokiej temperatury
      \item \textbf{Zasilanie:} zasilanie sieciowe, połączenie MagSafe o mocy 35W aby zagwarantować stały limit poboru mocy
        \item \textbf{Izolacja procesu:} Proces został przypisany do rdzeni wydajnościowych przy użyciu narzędzi dostępnych w systemie macOS (Quality of Service: User Interactive), poprzez nadanie mu wysokiego priorytetu za pomocą narzędzia nice za pomocą polecenia \textit{nice -n -20 $<scieżka do programu>$} 
        \item \textbf{Dokładny model urządzenia:} Apple A3241
      \end{itemize}
    \subsection{Instancje Badawcze}
    W Tab. \ref{tab:instancje} przedstawiono parametry wykorzystanych instancji badawczych do badania wszystkich
    algorytmów, a w Tab. \ref{tab:instancje2} przedstawiono dodatkowe zestawienie instancji dla algorytmu
    quicksort przeznaczone do badania wpływu typu danych na czas sortowania tablicy.
    Badanie każdej wielkości instancji powtórzono 100 razy aby ograniczyć wpływ potencjalnych
    błędów grubych podczas obliczania średniego czasu sortowania dla każdej instancji oraz osiągnąć statystycznie reprezentatywny wynik.
    \vspace{2 pt}
    \clearpage
    \begin{table}[h]
    \centering
    \caption{Wykaz i charakterystyka instancji badanych dla wszystkich algorytmów.}
    \label{tab:instancje}
    \begin{tabular}{|l|c|p{7cm}|c|}
      \hline
    \textbf{Lp.} & \textbf{Rozmiar ($N$)} & \textbf{Rodzaje rozkładu danych} & \textbf{Format danych} \\ \hline
    1 & 5 000  & losowy, posortowany (33\%, 66\%, 100\%), odwrotny & \texttt{int} \\ \hline
    2 & 10 000 & losowy, posortowany (33\%, 66\%, 100\%), odwrotny & \texttt{int} \\ \hline
    3 & 15 000 & losowy, posortowany (33\%, 66\%, 100\%), odwrotny & \texttt{int} \\ \hline
    4 & 20 000 & losowy, posortowany (33\%, 66\%, 100\%), odwrotny & \texttt{int} \\ \hline
    5 & 30 000 & losowy, posortowany (33\%, 66\%, 100\%), odwrotny & \texttt{int} \\ \hline
    6 & 40 000 & losowy, posortowany (33\%, 66\%, 100\%), odwrotny & \texttt{int} \\ \hline
    7 & 50 000 & losowy, posortowany (33\%, 66\%, 100\%), odwrotny & \texttt{int} \\ \hline
    \end{tabular}
    \small Źródło: opracowanie własne
  \end{table}
  \vspace{2 pt}
  \begin{table}[h]
    \centering
    \caption{Wykaz i charakterystyka dodatkowych instancji badanych dla sortowania szybkiego.}
    \label{tab:instancje2}
    \begin{tabular}{|l|c|p{7cm}|c|}
      \hline
    \textbf{Lp.} & \textbf{Rozmiar ($N$)} & \textbf{Rodzaje rozkładu danych} & \textbf{Format danych} \\ \hline
    1 & 5 000  & losowy, posortowany (33\%, 66\%, 100\%), odwrotny & \texttt{float, double, char} \\ \hline
    2 & 10 000 & losowy, posortowany (33\%, 66\%, 100\%), odwrotny & \texttt{float, double, char} \\ \hline
    3 & 15 000 & losowy, posortowany (33\%, 66\%, 100\%), odwrotny & \texttt{float, double, char} \\ \hline
    4 & 20 000 & losowy, posortowany (33\%, 66\%, 100\%), odwrotny & \texttt{float, double, char} \\ \hline
    5 & 30 000 & losowy, posortowany (33\%, 66\%, 100\%), odwrotny & \texttt{float, double, char} \\ \hline
    6 & 40 000 & losowy, posortowany (33\%, 66\%, 100\%), odwrotny & \texttt{float, double, char} \\ \hline
    7 & 50 000 & losowy, posortowany (33\%, 66\%, 100\%), odwrotny & \texttt{float, double, char} \\ \hline
    \end{tabular}
    \small Źródło: opracowanie własne
  \end{table} 
  \clearpage
  \section{Opis procedury badawczej}
    Procedura pomiarowa zakłada,
    że przypisanie procesowi wysokiego priorytetu systemowego zagwarantuje możliwie
    najstabilniejsze działanie programu i odizoluje go od wpływu innych procesów.
    Proces badawczy wykorzystywał bibliotekę \texttt{std::chrono},
    stosując zegar o wysokiej rozdzielczości \texttt{high\_resolution\_clock}.
    Wybór tego rozwiązania podyktowany był bardzo wysoką precyzją zegara,
    co jest kluczowe przy mierzeniu czasu niezwykle szybkich operacji sortowania.
    Pomiar obejmował wyłącznie właściwe operacje sortowania danych.
    Aby uniknąć przekłamania wyników przez czas generowania danych lub alokacji pamięci,
    algorytmy działały na wcześniej skopiowanych instancjach tablic.
    Weryfikacja poprawności sortowania,
    alokacja pamięci oraz kopiowanie danych były przeprowadzane poza mierzonym interwałem
    czasowym. Dla każdego rozmiaru instancji generowano i sortowano 100 różnych zestawów
    danych w celu uśrednienia wyników serii pomiarowej.
    Wynik końcowy prezentowany w analizie stanowi średnią arytmetyczną z wykonanych prób,
    co pozwoliło uzyskać statystycznie reprezentatywny wynik.
    W przeprowadzonym badaniu poprawność sortowania była monitorowana dla każdego pomiaru;
    wszystkie testy zakończyły się sukcesem,
    w związku z czym przypadek błędu w sortowaniu nie jest brany pod uwagę w dalszej analizie.
\section{Wyniki i ich analiza}

\begin{table}[H]
    \centering
    % 1. PODPIS NAD TABELĄ
    \caption{Zestawienie średnich czasów sortowania [ms] dla danych o rozkładzie losowym.}
    \label{tab:wyniki_losowe}
    
    \renewcommand{\arraystretch}{1.4} % Większy odstęp dla czytelności
    \small % Mniejsza czcionka, żeby tabela na pewno zmieściła się w marginesach
    
    % Definicja kolumn: m{szerokość} centruje w pionie, >{\centering...} w poziomie
    \begin{tabular}{|>{\centering\arraybackslash}m{1.5cm}|*{5}{>{\centering\arraybackslash}m{2.2cm}|}}
    \hline
    \textbf{Rozmiar} ($N$) & \textbf{Selection Sort} & \textbf{Insertion Sort} & \textbf{Heap Sort} & \textbf{Shell Sort} & \textbf{Quick Sort} \\ \hline
    
    % Wstawianie danych: używaj zapisu wykładniczego tam, gdzie liczby są bardzo małe lub duże
    5 000  & $4,12 \cdot 10^1$ & $2,45 \cdot 10^1$ & $1,12 \cdot 10^0$ & $2,10 \cdot 10^0$ & $8,45 \cdot 10^{-1}$ \\ \hline
    10 000 & $1,65 \cdot 10^2$ & $9,82 \cdot 10^1$ & $2,45 \cdot 10^0$ & $4,56 \cdot 10^0$ & $1,82 \cdot 10^0$ \\ \hline
    15 000 & $3,71 \cdot 10^2$ & $2,21 \cdot 10^2$ & $3,84 \cdot 10^0$ & $7,32 \cdot 10^0$ & $2,91 \cdot 10^0$ \\ \hline
    20 000 & $6,58 \cdot 10^2$ & $3,94 \cdot 10^2$ & $5,21 \cdot 10^0$ & $1,05 \cdot 10^1$ & $4,12 \cdot 10^0$ \\ \hline
    30 000 & $1,48 \cdot 10^3$ & $8,85 \cdot 10^2$ & $8,14 \cdot 10^0$ & $1,68 \cdot 10^1$ & $6,58 \cdot 10^0$ \\ \hline
    40 000 & $2,64 \cdot 10^3$ & $1,58 \cdot 10^3$ & $1,12 \cdot 10^1$ & $2,41 \cdot 10^1$ & $9,12 \cdot 10^0$ \\ \hline
    50 000 & $4,12 \cdot 10^3$ & $2,46 \cdot 10^3$ & $1,45 \cdot 10^1$ & $3,15 \cdot 10^1$ & $1,21 \cdot 10^1$ \\ \hline
    
    \end{tabular}
    
    \vspace{2pt}
    % 2. ŹRÓDŁO POD TABELĄ
    \small Źródło: opracowanie własne na podstawie przeprowadzonych badań.
\end{table}

    \clearpage
\printbibliography
\end{document}
